\chapter[Convergence: splitting, resonance, contract]{Phase 8: splitting with a boundary, periodic resonance, Mori--Zwanzig and the contract}\label{ch:cIV7}
\chapterbrief
{the notation of Chapter~\ref{ch:cIII3}. For B94: $I\subset\R$ open interval, $\mathbb T=\R/2\pi\Z$, $X=C^\infty(I\times\mathbb T)$. For B97: $X=L^2(0,L)$, $A_0=D\partial^2$ with Dirichlet conditions, $V=-c\partial_x$.}
{to give the proofs of the remaining convergence results: how periodic boundaries break the factor decomposition, the exact Mori--Zwanzig form for constant-coefficient systems, the error of a chain of steps, and the order of operator splitting with a Dirichlet boundary.}
{Dim.~B94--B97. They prove the corresponding statements of Chapter~\ref{PAP-ch:convergence} and of the extended algorithm (Chapter~\ref{PAP-ch:extalgo}).}

\begin{dimcard}{B94}{CONV-R1}{Periodic boundary: a finite-dimensional failure of the factor decomposition}{\proved}{}
Let $f,g\in R=\C[\tau,\xi]$ be coprime, $f_k(\tau)=f(\tau,ik)$, $g_k(\tau)=g(\tau,ik)$. Call $k\in\Z$ \emph{resonant} if $f_k\not\equiv0$, $g_k\not\equiv0$ and $\gcd(f_k,g_k)\ne1$. Then:
\begin{enumerate}[label=(\alph*)]
\item $\ker_X(fg)=\ker_Xf+\ker_Xg$ iff there is no resonant $k$;
\item the set of resonant $k$ is finite and $\dim\ker_X(fg)/(\ker_Xf+\ker_Xg)=\sum_{k\text{ resonant}}\deg\gcd(f_k,g_k)$;
\item on $I\times J$ with $J$ an interval (convex $\Omega$) the equality $\ker(fg)=\ker f+\ker g$ always holds (R2-T3, C12-T1$'$): the failure on $\mathbb T$ is a boundary phenomenon.
\end{enumerate}
\end{dimcard}

\begin{proof}
\step{1}{Bezout and finitely many resonances}
$f,g$ coprime in $\C[\xi][\tau]$ are coprime in $\C(\xi)[\tau]$ (Gauss), so $Af+Bg=r(\xi)$ with $0\ne r\in\C[\xi]$. Let $Z=\{k:r(ik)=0\}$, finite. For $k\notin Z$ a common root of $f_k,g_k$ would be a root of $r(ik)$: $k$ is not resonant. So the resonant $k$ lie in $Z$.

\step{2}{Reduction to Fourier modes}
For $v\in X$ put $\hat v_k(t)=(2\pi)^{-1}\int v\,e^{-ikx}\dd x$; then $(Q(\partial)v)^\wedge_k=Q(\partial_t,ik)\hat v_k$ and, by integration by parts in $x$, $\sup_K|\partial_t^j\hat v_k|\le C_{K,j,M}(1+|k|)^{-M}$; conversely a sequence with these bounds defines $\sum v_ke^{ikx}\in X$. For non-zero $p,q\in\C[\tau]$ on $C^\infty(I)$: $\dim\ker p(\partial_t)=\deg p$ and $\ker p\cap\ker q=\ker\gcd(p,q)$, so $\mathrm{codim}(\ker p+\ker q)$ in $\ker(pq)$ is $\deg\gcd(p,q)$.

\step{3}{Sufficiency: an explicit splitting}
Let $u\in\ker_X(fg)$, so $\hat u_k\in\ker(f_kg_k)$. For $k\notin Z$ put $v_k=B_kg_k\hat u_k/r(ik)$, $w_k=A_kf_k\hat u_k/r(ik)$: then $v_k+w_k=\hat u_k$, $f_kv_k=0$, $g_kw_k=0$. For the finitely many non-resonant $k\in Z$ split by the ODE facts. Regularity: $|r(ik)|\ge c>0$ for large $|k|$ (a non-zero polynomial), and $A_k,B_k,f_k,g_k$ are differential operators of fixed order, so the decay in $k$ is preserved and $\sum v_ke^{ikx},\sum w_ke^{ikx}\in X$.

\step{4}{Necessity and the dimension}
If $k$ is resonant take $y\in\ker(f_kg_k)\setminus(\ker f_k+\ker g_k)$ and $u=y(t)e^{ikx}$. A splitting $u=v+w$ would give $y=\hat v_k+\hat w_k$ with $f_k\hat v_k=0$, $g_k\hat w_k=0$: contradiction. More precisely, $u\mapsto(\hat u_k)_{k\text{ resonant}}$ induces an isomorphism $\ker(fg)/(\ker f+\ker g)\cong\bigoplus_{k\text{ res.}}\ker(f_kg_k)/(\ker f_k+\ker g_k)$, onto by finite sums, with kernel $\ker f+\ker g$ by Step~3 on non-resonant modes.

\whyworks{on a circle every mode is an ODE in $t$, and a common root of the two factors at a frequency $ik$ is a genuine overlap that Bezout in one variable cannot split.}
\end{proof}
\uses{A13, A32 (Bezout).}

\begin{example}{Resonance at $k=0$}
$f=\tau-\xi^2$, $g=\tau+\xi$: $f-g=-\xi(\xi+1)$, the only resonant mode is $k=0$ ($f_0=g_0=\tau$), the codimension is $1$, and the witness is $u=t$ (on $\R$, $t=x-(x-t)$). For $f=\tau-\xi^2$, $g=\tau+\xi+1$ the common roots satisfy $\xi^2+\xi+1=0$, outside $i\Z$: the decomposition always exists. For $f=\xi$, $g=\tau$ the line $\xi=0$ lies in $V(f)$, so $(0,0)$ is not resonant and $\ker(\tau\xi)=\{A(t)+B(x)\}=\ker\xi+\ker\tau$.
\end{example}

\begin{falsebox}{The criterion without the clause on contained lines}{\falsest}
The criterion ``no resonance $V(f)\cap V(g)\cap(\C\times i\Z)$'' of the route RT-R2a is correct only up to lines contained in $V(f)$ or $V(g)$; without that clause the example $f=\xi$, $g=\tau$ refutes it (CONV-R1-F).
\end{falsebox}

\begin{dimcard}{B95}{CONV-MZ}{Mori--Zwanzig with constant coefficients}{\proved}{}
\begin{enumerate}[label=(MZ\arabic*)]
\item If $D\in\mathrm{Mat}_k(R_n)$, $\det D\ne0$, $\mathcal A$ an $R_n$-module and $DU=0$, then $(\det D)U_i=0$ for every $i$: each component of a square non-degenerate system has a \emph{local} law of order at most $\deg\det D$.
\item For $(\hat u,\hat v)'=A(i\kappa)(\hat u,\hat v)$, $A=\begin{pmatrix}a_{11}&a_{12}\\a_{21}&a_{22}\end{pmatrix}$, $a_{ij}\in\C[\xi]$:
\[
\hat u'(t)=a_{11}\hat u(t)+a_{12}\!\int_0^t\!e^{(t-s)a_{22}}a_{21}\hat u(s)\dd s+a_{12}e^{ta_{22}}\hat v(0)
\]
(Markov term, memory, ``noise'' from unresolved data), and $(\partial_t-a_{22})(\partial_t-a_{11})\hat u=a_{12}a_{21}\hat u$, i.e.\ $\det(\tau-A(\xi))(\partial)u=0$.
\item The memory kernel $K(t,\xi)=a_{12}e^{ta_{22}}a_{21}$ vanishes iff $a_{12}a_{21}=0$; $[\mathcal P,A]=0$ with $\mathcal P=\mathrm{diag}(1,0)$ iff $a_{12}=a_{21}=0$. The criterion ``$E_{\rm int}=0\iff[\mathcal P,L]=0$'', read as ``memory zero $\iff[\mathcal P,A]=0$'', is \textbf{false}, and so is the reading ``memory and noise zero $\iff[\mathcal P,A]=0$''.
\end{enumerate}
\end{dimcard}

\begin{proof}
\step{1}{(MZ1)}
$\mathrm{adj}(D)D=\det(D)I$, so $\det(D)U=\mathrm{adj}(D)(DU)=0$.

\step{2}{(MZ2) Duhamel}
$\hat v(t)=e^{ta_{22}}\hat v(0)+\int_0^te^{(t-s)a_{22}}a_{21}\hat u(s)\dd s$; substitute in $\hat u'=a_{11}\hat u+a_{12}\hat v$. Applying $\partial_t-a_{22}$ to $\hat u'-a_{11}\hat u=a_{12}\hat v$ gives $a_{12}(\hat v'-a_{22}\hat v)=a_{12}a_{21}\hat u$.

\step{3}{(MZ3) and the counterexamples}
$\C[\xi]$ is a domain and $e^{ta_{22}}\ne0$, so $K\equiv0\iff a_{12}a_{21}=0$; $\mathcal PA-A\mathcal P=\begin{pmatrix}0&a_{12}\\-a_{21}&0\end{pmatrix}$. With $a_{21}=0$, $a_{12}=1$ the memory is zero but $[\mathcal P,A]\ne0$. With $a_{12}=0$, $a_{21}=1$ the component $\hat u$ is exactly closed ($\hat u'=a_{11}\hat u$, memory and noise zero) but $[\mathcal P,A]=\begin{pmatrix}0&0\\-1&0\end{pmatrix}\ne0$.

\whyworks{with constant coefficients, eliminating a variable gives back a local law (the determinant), and ``memory'' is the same information written as an integral.}
\end{proof}
\uses{none.}

\begin{caution}[title=Reading]
For under-determined systems the eliminated law may be missing ($u_t=v_x$: every $u$ is a component of a solution on convex $\Omega$, $v=\int u_t\dd x$), consistent with the class $\mathcal N_{\mathcal K}$ of B91. With variable coefficients the converse ``closed local law $\iff$ finite elimination'' is open (CONV-Q11).
\end{caution}

\begin{dimcard}{B96}{CONV-T7}{The error of a chain of steps}{\proved}{}
Let $X_0,\dots,X_k$ be metric spaces, $F_i:X_{i-1}\to X_i$ the exact maps and $\tilde F_i$ the computed ones, $\tilde x_0=x_0$, $\tilde x_i=\tilde F_i(\tilde x_{i-1})$, $x_i=F_i(x_{i-1})$. If $d(\tilde F_i(\tilde x_{i-1}),F_i(\tilde x_{i-1}))\le\varepsilon_i$ and each $F_j$ is $L_j$-Lipschitz on a set containing the points that occur below, then
\[
d(\tilde x_k,x_k)\le\sum_{i=1}^k\Big(\prod_{j>i}L_j\Big)\varepsilon_i .
\]
\end{dimcard}

\begin{proof}
\step{1}{Telescoping}
Put $y_i=F_k\circ\dots\circ F_{i+1}(\tilde x_i)$, so $y_k=\tilde x_k$ and $y_0=x_k$. Then $d(\tilde x_k,x_k)\le\sum_{i=1}^kd(y_i,y_{i-1})$.

\step{2}{Each term}
$y_{i-1}=F_k\circ\dots\circ F_{i+1}(F_i(\tilde x_{i-1}))$ and $y_i=F_k\circ\dots\circ F_{i+1}(\tilde F_i(\tilde x_{i-1}))$, so $d(y_i,y_{i-1})\le(\prod_{j>i}L_j)\varepsilon_i$.

\whyworks{an error made at step $i$ is amplified only by the Lipschitz constants of the steps that follow.}
\end{proof}
\uses{none.}

\begin{caution}[title=Use in the contract]
The guaranteed chain contains only steps of class (A) (finite or exact: elimination, Wei--Norman, $E_{\rm int}$ in the class $\mathfrak B$, commuting multipliers, MZ1) and (B) (explicit $\varepsilon_i$: Dyson--Phillips with a bounded perturbation, Magnus only conditionally). A step of class (C) has no $\varepsilon_i$ and is outside the guarantee; a step of class (Lim) is exact from $K^*$ on, but $K^*$ is not certifiable, so the chain guarantee holds ``under the event $K\ge K^*$'', which must be declared to the user. As a theorem the proposition is trivial; its value is as a specification of Step~7.
\end{caution}

\begin{dimcard}{B97}{CONV-T8}{Splitting with a Dirichlet boundary: orders in $L^2$}{\proved}{}
Let $e_n=\sqrt{2/L}\sin(\mu_nx)$, $\mu_n=n\pi/L$, $\lambda_n=-\kappa n^2$, $\kappa=D\pi^2/L^2$, $V_{mn}=\langle e_m,Ve_n\rangle=-4cmn/(L(m^2-n^2))$ for $m+n$ odd, $0$ otherwise; $S^{(1)}(t)=\int_0^tS_0(t-s)VS_0(s)\dd s=\partial_\varepsilon S_\varepsilon|_{\varepsilon=0}$; $h>0$ the step, $n_h=T/h$. Let $G_h^{\rm Lie}$, $G_h^{\rm Str}$ be the first-order-in-$\varepsilon$ global errors of the Lie and Strang splittings for modal data $u_0=\sum_{n\le N}c_ne_n$.
\begin{enumerate}[label=(\alph*)]
\item (closed form) with $x_{mn}=\kappa(m^2-n^2)h$, $\varphi_L(x)=1/x-1/(e^x-1)$, $\varphi_S(x)=1/x-\frac12\coth(x/2)$:
\[
\langle e_m,G_h^{\rm Lie}\rangle=\sum_nc_nV_{mn}(e^{\lambda_nT}-e^{\lambda_mT})\,h\,\varphi_L(x_{mn}),\qquad \langle e_m,G_h^{\rm Str}\rangle=\sum_nc_nV_{mn}(e^{\lambda_nT}-e^{\lambda_mT})\,h\,\varphi_S(x_{mn}).
\]
\item (Lie: global order $1$) $\|G_h^{\rm Lie}\|\le h\sqrt N(\sum|c_n|^2c^2\mu_n^2)^{1/2}$ and, for $u_0=e_1$, $\|G_h^{\rm Lie}\|/h\to C_L>0$.
\item (Strang: order reduced from $2$ to $5/4$ in $L^2$) $\|G_h^{\rm Str}\|\le Ch^{5/4}$ for $0<h\le1$, and for $u_0=e_1$ and $h\le1/(36\kappa)$ a lower bound $\ge c_*h^{5/4}$.
\end{enumerate}
\end{dimcard}

\begin{proof}
\step{1}{(a) Explicit sums}
By linearity take $u_0=e_n$; $\langle e_m,S_0(r)VS_0(s)e_n\rangle=e^{\lambda_mr}V_{mn}e^{\lambda_ns}$ and $\langle e_m,S^{(1)}(T)e_n\rangle=V_{mn}(e^{\lambda_nT}-e^{\lambda_mT})/(\lambda_n-\lambda_m)$. With $a=e^{\lambda_nh}$, $b=e^{\lambda_mh}$ the Lie sum is $V_{mn}hb(a^{n_h}-b^{n_h})/(a-b)$, and $b/(a-b)=1/(e^x-1)$, $1/\Delta=h/x$ with $x=\Delta h$, giving $\varphi_L$; the Strang sum gives $\varphi_S$ in the same way.

\step{2}{Properties of $\varphi_L,\varphi_S$}
$\varphi_L\in(0,1)$ for $x\ne0$ and $\varphi_L(x)\to\frac12$ as $x\to0$. $\varphi_S$ is odd, and for $y>0$, $\varphi_S(y)=-\frac12\mathcal L(y/2)$ with the Langevin function $\mathcal L(z)=\coth z-1/z$, increasing with $0<\mathcal L(z)<\min(z/3,1)$; so $|\varphi_S(x)|\le\min(|x|/12,\frac12)$ and $|\varphi_S|$ increases in $|x|$.

\step{3}{(b) Lie}
$|e^{\lambda_nT}-e^{\lambda_mT}|\le1$ and $|\varphi_L|\le1$ give $|\langle e_m,G^{\rm Lie}\rangle|\le h\sum_n|c_n||V_{mn}|$; Cauchy--Schwarz and Parseval ($\|Ve_n\|=|c|\mu_n$) give the bound. For $u_0=e_1$ and fixed $m$, $x_{m1}\to0$ so $\varphi_L\to\frac12$; dominated convergence gives $\|G^{\rm Lie}\|^2/h^2\to C_L^2=\frac14\sum_{m\text{ even}}V_{m1}^2(e^{\lambda_1T}-e^{\lambda_mT})^2>0$ since $V_{21}\ne0$.

\step{4}{(c) Strang: the upper bound}
For $m\le2N$, $|\varphi_S(x_{mn})|\le\kappa4N^2h/12$ and $|V_{mn}|\le8|c|N^2/L$, so the contribution is $O(h^4)$ in the squared norm. For $m>2N$, $m^2-n^2\ge\frac34m^2$ gives $|\langle e_m,G^{\rm Str}\rangle|\le\frac{16|c|A}{3L}m^{-1}h\min(\kappa m^2h/12,\frac12)$; splitting at $M=\lfloor h^{-1/2}\rfloor$, $\sum_{m\le M}\lesssim\kappa^2h^{1/2}$ and $\sum_{m>M}\lesssim h^{1/2}$, so $\|G^{\rm Str}\|^2\le Ch^{5/2}$. (The lower bound uses the window of even $m$ with $1\le x_{m1}\le4$ and is \emph{condensed}: package \S6.)

\whyworks{the transport step does not preserve the Dirichlet condition ($Ve_1\notin H_0^1$), producing a boundary layer on the modes $m\sim h^{-1/2}$; the parabolic smoothing absorbs it for Lie, but for Strang it lowers the order from $2$ to $5/4$.}
\end{proof}
\uses{none.}

\begin{caution}[title=Scope and priority]
First order in $\varepsilon$ only, modal data, transport step truncated at first order. In $L^\infty$ near $x=0$ the observed Strang slope is $1.004\to1.001$ (order $1$, checked on an example, not proved). The mechanism is that of Einkemmer--Ostermann (2015) for diffusion--reaction splitting with non-compatible boundary conditions; the norm in which Dang--Einkemmer--Ostermann (2025) state their order is not verified (CONV-Q6$'$).
\end{caution}
